% GENERATED FILE. Edit canonical lessons, then run npm run generate:books.
% canonical-source-hash: fnv1a64:af24a7bda1b8daa6
% canonical-lessons: KA-C135-galla, KA-C135-kenne, KA-C135-monakalu, KA-C135-monakai, KA-C135-manikattu, KA-R135-first-pass-food-and-nature, KA-R135-second-pass-nature-and-animals

\chapter{Chin, Cheek, Knee, Elbow, Wrist}
\label{ch:ka-chin-cheek-knee-elbow-wrist}
\hlchaptermodality{135}

\begin{chapteropening}
\textbf{By the end of this chapter:} \emph{I can say the chin, a cheek, a knee, an elbow and a wrist.}

Complete the last lesson of chapter 135: i can say the chin, a cheek, a knee, an elbow and a wrist.
\end{chapteropening}

\section[\texorpdfstring{\kn{ಗಲ್ಲ} (galla)}{ಗಲ್ಲ (galla)}]{\kn{ಗಲ್ಲ} (galla) — the chin}

\label{lesson:KA-C135-galla}



\begin{quote}
Before the new one: say the Kannada for the face, then the Kannada for the tongue.
\end{quote}

\subsection*{You'll want to know: \kn{ಗಲ್ಲ}}
\textbf{\kn{ಗಲ್ಲ}} — \emph{galla} — \textquotedblleft{}the chin\textquotedblright{}.

\begin{grammarlens}[title={What you've built}]
One more word for this chapter.
\end{grammarlens}

\subsection*{Your turn}
\begin{itemize}
\raggedright
  \item \emph{Say it:} \emph{galla}
  \item \emph{Say it:} \emph{galla}, once more
  \item \emph{Say it:} say \emph{galla}
  \item \emph{Recall:} say the Kannada for the face, then the Kannada for the tongue, then say \emph{galla} again
\end{itemize}

\begin{tcolorbox}[breakable,colback=teal!4,colframe=teal!35!black,title={Before you move on}]
What does \kn{ಗಲ್ಲ} mean? (\textquotedblleft{}The chin\textquotedblright{}.) Say it once more.
\end{tcolorbox}
\section[\texorpdfstring{\kn{ಕೆನ್ನೆ} (kenne)}{ಕೆನ್ನೆ (kenne)}]{\kn{ಕೆನ್ನೆ} (kenne) — a cheek}

\label{lesson:KA-C135-kenne}



\begin{quote}
Before the new one: say the Kannada for the tongue, then the Kannada for the chin.
\end{quote}

\subsection*{You'll want to know: \kn{ಕೆನ್ನೆ}}
\textbf{\kn{ಕೆನ್ನೆ}} — \emph{kenne} — \textquotedblleft{}a cheek\textquotedblright{}.

\begin{grammarlens}[title={What you've built}]
One more word for this chapter.
\end{grammarlens}

\subsection*{Your turn}
\begin{itemize}
\raggedright
  \item \emph{Say it:} \emph{kenne}
  \item \emph{Say it:} \emph{kenne}, once more
  \item \emph{Say it:} say \emph{kenne}
  \item \emph{Recall:} say the Kannada for the tongue, then the Kannada for the chin, then say \emph{kenne} again
\end{itemize}

\begin{tcolorbox}[breakable,colback=teal!4,colframe=teal!35!black,title={Before you move on}]
What does \kn{ಕೆನ್ನೆ} mean? (\textquotedblleft{}A cheek\textquotedblright{}.) Say it once more.
\end{tcolorbox}
\section[\texorpdfstring{\kn{ಮೊಣಕಾಲು} (moṇakālu)}{ಮೊಣಕಾಲು (moṇakālu)}]{\kn{ಮೊಣಕಾಲು} (moṇakālu) — a knee}

\label{lesson:KA-C135-monakalu}



\begin{quote}
Before the new one: say the Kannada for the chin, then the Kannada for a cheek.
\end{quote}

\subsection*{You'll want to know: \kn{ಮೊಣಕಾಲು}}
\textbf{\kn{ಮೊಣಕಾಲು}} — \emph{moṇakālu} — \textquotedblleft{}a knee\textquotedblright{}.

\begin{grammarlens}[title={What you've built}]
One more word for this chapter.
\end{grammarlens}

\subsection*{Your turn}
\begin{itemize}
\raggedright
  \item \emph{Say it:} \emph{moṇakālu}
  \item \emph{Say it:} \emph{moṇakālu}, once more
  \item \emph{Say it:} say \emph{moṇakālu}
  \item \emph{Recall:} say the Kannada for the chin, then the Kannada for a cheek, then say \emph{moṇakālu} again
\end{itemize}

\begin{tcolorbox}[breakable,colback=teal!4,colframe=teal!35!black,title={Before you move on}]
What does \kn{ಮೊಣಕಾಲು} mean? (\textquotedblleft{}A knee\textquotedblright{}.) Say it once more.
\end{tcolorbox}
\section[\texorpdfstring{\kn{ಮೊಣಕೈ} (moṇakai)}{ಮೊಣಕೈ (moṇakai)}]{\kn{ಮೊಣಕೈ} (moṇakai) — an elbow}

\label{lesson:KA-C135-monakai}



\begin{quote}
Before the new one: say the Kannada for a cheek, then the Kannada for a knee.
\end{quote}

\subsection*{You'll want to know: \kn{ಮೊಣಕೈ}}
\textbf{\kn{ಮೊಣಕೈ}} — \emph{moṇakai} — \textquotedblleft{}an elbow\textquotedblright{}.

\begin{grammarlens}[title={What you've built}]
One more word for this chapter.
\end{grammarlens}

\subsection*{Your turn}
\begin{itemize}
\raggedright
  \item \emph{Say it:} \emph{moṇakai}
  \item \emph{Say it:} \emph{moṇakai}, once more
  \item \emph{Say it:} say \emph{moṇakai}
  \item \emph{Recall:} say the Kannada for a cheek, then the Kannada for a knee, then say \emph{moṇakai} again
\end{itemize}

\begin{tcolorbox}[breakable,colback=teal!4,colframe=teal!35!black,title={Before you move on}]
What does \kn{ಮೊಣಕೈ} mean? (\textquotedblleft{}An elbow\textquotedblright{}.) Say it once more.
\end{tcolorbox}
\section[\texorpdfstring{\kn{ಮಣಿಕಟ್ಟು} (maṇikaṭṭu)}{ಮಣಿಕಟ್ಟು (maṇikaṭṭu)}]{\kn{ಮಣಿಕಟ್ಟು} (maṇikaṭṭu) — a wrist}

\label{lesson:KA-C135-manikattu}



\begin{quote}
Before the new one: say the Kannada for a knee, then the Kannada for an elbow.
\end{quote}

\subsection*{You'll want to know: \kn{ಮಣಿಕಟ್ಟು}}
\textbf{\kn{ಮಣಿಕಟ್ಟು}} — \emph{maṇikaṭṭu} — \textquotedblleft{}a wrist\textquotedblright{}.

\begin{grammarlens}[title={What you've built}]
One more word for this chapter.
\end{grammarlens}

\subsection*{Your turn}
\begin{itemize}
\raggedright
  \item \emph{Say it:} \emph{maṇikaṭṭu}
  \item \emph{Say it:} \emph{maṇikaṭṭu}, once more
  \item \emph{Say it:} say \emph{maṇikaṭṭu}
  \item \emph{Recall:} say the Kannada for a knee, then the Kannada for an elbow, then say \emph{maṇikaṭṭu} again
\end{itemize}

\begin{tcolorbox}[breakable,colback=teal!4,colframe=teal!35!black,title={Before you move on}]
What does \kn{ಮಣಿಕಟ್ಟು} mean? (\textquotedblleft{}A wrist\textquotedblright{}.) Say it once more.
\end{tcolorbox}
\section[(dialogue)]{Food and nature — a first pass}

\label{lesson:KA-R135-first-pass-food-and-nature}



\begin{quote}
Say the words for an elbow and a wrist.
\end{quote}

\subsection*{Your turn}
Say these aloud:

\begin{itemize}
\raggedright
  \item the sea, a forest, soil, a stone, thunder
  \item lightning, a storm, a wave, a waterfall, an island
  \item the sun, a horse, an elephant, a tiger, a lion
  \item a bear, a monkey, a rabbit, a deer, a fox
\end{itemize}

\begin{tcolorbox}[breakable,colback=teal!4,colframe=teal!35!black,title={Before you move on}]
The sea? (\textbf{\kn{ಸಮುದ್ರ}}, \emph{samudra}.) A forest? (\textbf{\kn{ಕಾಡು}}, \emph{kāḍu}.) Soil? (\textbf{\kn{ಮಣ್ಣು}}, \emph{maṇṇu}.) A stone? (\textbf{\kn{ಕಲ್ಲು}}, \emph{kallu}.) Thunder? (\textbf{\kn{ಗುಡುಗು}}, \emph{guḍugu}.) Lightning? (\textbf{\kn{ಮಿಂಚು}}, \emph{miṁcu}.) A storm? (\textbf{\kn{ಬಿರುಗಾಳಿ}}, \emph{birugāḷi}.) A wave? (\textbf{\kn{ಅಲೆ}}, \emph{ale}.) A waterfall? (\textbf{\kn{ಜಲಪಾತ}}, \emph{jalapāta}.) An island? (\textbf{\kn{ದ್ವೀಪ}}, \emph{dvīpa}.) The sun? (\textbf{\kn{ಸೂರ್ಯ}}, \emph{sūrya}.) A horse? (\textbf{\kn{ಕುದುರೆ}}, \emph{kudure}.) An elephant? (\textbf{\kn{ಆನೆ}}, \emph{āne}.) A tiger? (\textbf{\kn{ಹುಲಿ}}, \emph{huli}.) A lion? (\textbf{\kn{ಸಿಂಹ}}, \emph{siṁha}.) A bear? (\textbf{\kn{ಕರಡಿ}}, \emph{karaḍi}.) A monkey? (\textbf{\kn{ಕೋತಿ}}, \emph{kōti}.) A rabbit? (\textbf{\kn{ಮೊಲ}}, \emph{mola}.) A deer? (\textbf{\kn{ಜಿಂಕೆ}}, \emph{jiṅke}.) A fox? (\textbf{\kn{ನರಿ}}, \emph{nari}.)
\end{tcolorbox}
\section[(dialogue)]{Nature and animals — a second pass}

\label{lesson:KA-R135-second-pass-nature-and-animals}



\begin{quote}
Say the words for a snake and a tortoise.
\end{quote}

\subsection*{Your turn}
Say these aloud:

\begin{itemize}
\raggedright
  \item a snake, a tortoise, a parrot, a pigeon, an owl
  \item a peacock, a butterfly, an ant, a mosquito, a spider
  \item a frog, a duck, a pig, a donkey, a camel
  \item a rat, a buffalo, a sheep, the face, the tongue
  \item the chin, a cheek, a knee, an elbow, a wrist
\end{itemize}

\begin{tcolorbox}[breakable,colback=teal!4,colframe=teal!35!black,title={Before you move on}]
A snake? (\textbf{\kn{ಹಾವು}}, \emph{hāvu}.) A tortoise? (\textbf{\kn{ಆಮೆ}}, \emph{āme}.) A parrot? (\textbf{\kn{ಗಿಳಿ}}, \emph{giḷi}.) A pigeon? (\textbf{\kn{ಪಾರಿವಾಳ}}, \emph{pārivāḷa}.) An owl? (\textbf{\kn{ಗೂಬೆ}}, \emph{gūbe}.) A peacock? (\textbf{\kn{ನವಿಲು}}, \emph{navilu}.) A butterfly? (\textbf{\kn{ಚಿಟ್ಟೆ}}, \emph{ciṭṭe}.) An ant? (\textbf{\kn{ಇರುವೆ}}, \emph{iruve}.) A mosquito? (\textbf{\kn{ಸೊಳ್ಳೆ}}, \emph{soḷḷe}.) A spider? (\textbf{\kn{ಜೇಡ}}, \emph{jēḍa}.) A frog? (\textbf{\kn{ಕಪ್ಪೆ}}, \emph{kappe}.) A duck? (\textbf{\kn{ಬಾತುಕೋಳಿ}}, \emph{bātukōḷi}.) A pig? (\textbf{\kn{ಹಂದಿ}}, \emph{handi}.) A donkey? (\textbf{\kn{ಕತ್ತೆ}}, \emph{katte}.) A camel? (\textbf{\kn{ಒಂಟೆ}}, \emph{oṇṭe}.) A rat? (\textbf{\kn{ಇಲಿ}}, \emph{ili}.) A buffalo? (\textbf{\kn{ಎಮ್ಮೆ}}, \emph{emme}.) A sheep? (\textbf{\kn{ಕುರಿ}}, \emph{kuri}.) The face? (\textbf{\kn{ಮುಖ}}, \emph{mukha}.) The tongue? (\textbf{\kn{ನಾಲಿಗೆ}}, \emph{nālige}.) The chin? (\textbf{\kn{ಗಲ್ಲ}}, \emph{galla}.) A cheek? (\textbf{\kn{ಕೆನ್ನೆ}}, \emph{kenne}.) A knee? (\textbf{\kn{ಮೊಣಕಾಲು}}, \emph{moṇakālu}.) An elbow? (\textbf{\kn{ಮೊಣಕೈ}}, \emph{moṇakai}.) A wrist? (\textbf{\kn{ಮಣಿಕಟ್ಟು}}, \emph{maṇikaṭṭu}.)
\end{tcolorbox}
